\documentclass[12pt, a4paper]{ctexart}
\usepackage[margin=2cm]{geometry}
\usepackage{libertine}

\usepackage{titlesec}
\usepackage{zhnumber}
\titleformat*{\section}{\Large\bfseries\raggedright}
\renewcommand{\thesection}{\zhnum{section}}
\renewcommand{\thesubsection}{\arabic{subsection}}

\setcounter{secnumdepth}{4}
\renewcommand{\theparagraph}{\thesubsubsection.\arabic{paragraph}}
\renewcommand{\paragraph}[1]{
    \refstepcounter{paragraph}
    {\bfseries\theparagraph\quad#1}\par
    \vspace{2pt}
    \noindent
}

\usepackage{enumitem}
\setlist[enumerate]{itemsep=2pt, parsep=0pt, partopsep=0pt, topsep=2pt}
\setlist[itemize]{itemsep=2pt, parsep=0pt, partopsep=0pt, topsep=2pt}
\linespread{1.2}

\usepackage{minted}
\setminted{
    breaklines=true, escapeinside=||, fontsize=\small, frame=lines,
    mathescape=$$, numbers=none, style=bw, tabsize=4
}

\usepackage{graphicx}
\usepackage{siunitx}

\begin{document}

\pagestyle{plain}
\thispagestyle{empty}

\noindent
\begin{tabular*}{\textwidth}{l @{\extracolsep{\fill}} r @{\extracolsep{6pt}} l}
    \LARGE{\textbf{实验报告}} & 数据结构 & \textit{Data Structures} \\
\end{tabular*}\\\\
\begin{tabular*}{\textwidth}{l l}
    \textbf{报告标题: } & \textbf{拓扑排序} \\
    \textbf{学号: } & 19240212 \\
    \textbf{姓名: } & 华博文 \\
    \textbf{日期: } & 2025 年 12 月 9 日 \\
\end{tabular*}\\
\rule[2ex]{\textwidth}{2pt}

\section{实验环境}

\subsection{操作系统}

Ubuntu 24.04.3 LTS x86\_64，Linux 6.14.0-36-generic，AMD Ryzen 7 7700 (16) @ 5.3GHz，DDR5。

\subsection{编程工具}

NeoVim v0.12.0-dev，gcc 13.3.0。

\subsection{其他工具}

\LaTeX{} (\TeX\textit{studio})。

\section{实验目的}

\begin{enumerate}
    \item 理解有向无环图（DAG）的概念与特性；
    \item 掌握拓扑排序算法的原理与实现方法；
    \item 学习利用拓扑排序检测有向图中是否存在回路；
    \item 了解拓扑排序在任务调度、依赖分析等实际问题中的应用。
\end{enumerate}

\section{实验内容及其完成情况}

\begin{enumerate}
    \item 实现有向图的拓扑排序：使用邻接表存储图，从文件读取输入；
    \item 计算并维护每个顶点的入度，使用队列管理入度为 $0$ 的顶点，输出拓扑序列；
    \item 若处理顶点数小于总顶点数，判定存在环并报告；
    \item 输入输出格式清晰，能检测异常输入并给出提示。
\end{enumerate}

\subsection{数据结构}

\begin{itemize}
    \item 邻接表：\mintinline{c++}`std::vector<std::vector<std::size_t>> adj`；
    \item 入度表：\mintinline{c++}`std::vector<std::size_t> indegree`；
    \item 就绪队列：\mintinline{c++}`std::queue<std::size_t>`（FIFO，无需最小编号优先）。
\end{itemize}

\subsection{程序实现}

输入的首行是 \mintinline{text}`n m`，接着输入 \mintinline{text}`m` 行边 \mintinline{text}`u v`（1-based 顶点编号，表示有向边 $u\to v$）。支持传入 \mintinline{text}`-` 从标准输入读取。

图存储与入度统计：

\begin{minted}{c++}
struct Graph {
    std::vector<std::vector<std::size_t>> adj;
    std::vector<std::size_t> indegree;
};

Graph parse_graph(std::istream &in) {
    std::size_t n, m;
    if (!(in >> n >> m))
        throw std::runtime_error("invalid header (expected n m)");
    if (n == 0)
        throw std::runtime_error("n must be > 0");

    Graph g;
    g.adj.assign(n, {});
    g.indegree.assign(n, 0);

    for (std::size_t i = 0; i < m; ++i) {
        long u, v;
        if (!(in >> u >> v))
            throw std::runtime_error("invalid edge line");
        if (u < 1 || v < 1 || u > static_cast<long>(n) || v > static_cast<long>(n)) {
            throw std::runtime_error("vertex id out of range; use 1..n");
        }
        const std::size_t from = static_cast<std::size_t>(u - 1);
        const std::size_t to = static_cast<std::size_t>(v - 1);
        g.adj[from].push_back(to);
        ++g.indegree[to];
    }

    return g;
}

Graph load_graph_from(const std::string &path) {
    if (path == "-")
        return parse_graph(std::cin);
    std::ifstream fin(path);
    if (!fin)
        throw std::runtime_error("cannot open input file: " + path);
    return parse_graph(fin);
}
\end{minted}

拓扑排序使用 Kahn 算法实现：

\begin{enumerate}
    \item 统计各顶点入度，将入度为 $0$ 的顶点入队；
    \item 循环弹出队首顶点 \mintinline{text}`u`，加入拓扑序，对其出边 \mintinline{text}`u -> v` 令 \mintinline{c++}`indeg[v]--`，若降为 $0$ 则入队；
    \item 结束后，若已输出顶点数 $<n$，则存在环。
\end{enumerate}

\begin{minted}{c++}
std::optional<std::vector<std::size_t>> topo_sort(const Graph &g, std::size_t &processed) {
    std::vector<std::size_t> indeg = g.indegree;
    std::queue<std::size_t> ready;
    for (std::size_t i = 0; i < indeg.size(); ++i) {
        if (indeg[i] == 0)
            ready.push(i);
    }

    std::vector<std::size_t> order;
    order.reserve(indeg.size());

    while (!ready.empty()) {
        std::size_t u = ready.front();
        ready.pop();
        order.push_back(u);
        for (std::size_t v : g.adj[u]) {
            if (--indeg[v] == 0)
                ready.push(v);
        }
    }

    processed = order.size();
    if (processed != g.adj.size())
        return std::nullopt;
    return order;
}
\end{minted}

输出与错误处理：

\begin{minted}{c++}
void print_topo_order(const std::vector<std::size_t> &order) {
    std::cout << "Topological order: ";
    for (std::size_t i = 0; i < order.size(); ++i) {
        if (i)
            std::cout << ' ';
        std::cout << (order[i] + 1);
    }
    std::cout << "\n";
}

int main(int argc, char **argv) {
    if (argc < 2) {
        std::cerr << "Usage: " << argv[0] << " <graph.txt|->\n";
        std::cerr << "Input format: n m followed by m lines of edges u v." << std::endl;
        return 1;
    }

    try {
        Graph g = load_graph_from(argv[1]);
        std::size_t processed = 0;
        auto order = topo_sort(g, processed);
        if (!order) {
            std::cout << "Cycle detected (processed " << processed << "/" << g.adj.size() << ").\n";
            return 0;
        }
        print_topo_order(*order);
    } catch (const std::exception &e) {
        std::cerr << "Error: " << e.what() << "\n";
        return 1;
    }

    return 0;
}
\end{minted}

\subsection{编译与运行示例}

在 \mintinline{text}`src` 目录下编译：

\begin{minted}{bash}
g++ -std=c++17 main.cpp -o topo
\end{minted}

运行时，图可以从文件读取：

\begin{minted}{bash}
./topo ../res/dag.txt
\end{minted}

也可以从标准输入读取（\mintinline{text}`-` 表示 stdin）：

\begin{minted}{bash}
cat ../res/dag.txt | ./topo -
\end{minted}

\subsection{样例输入输出}

无环示例（\mintinline{bash}`res/dag.txt`）：

\inputminted{text}{../res/dag.txt}

输出 \mintinline{text}`Topological order: 1 2 3 4 5 6`。

有环示例（\mintinline{bash}`res/cycle.txt`）：

\inputminted{text}{../res/cycle.txt}

输出 \mintinline{text}`Cycle detected (processed 0/3)`。

\subsection{复杂度分析}

\begin{itemize}
    \item 时间复杂度：$\mathcal{O}\left(n+m\right)$，每条边与顶点各处理一次；
    \item 空间复杂度：$\mathcal{O}\left(n+m\right)$，邻接表与入度数组、队列存储。
\end{itemize}

\section{结论}

实现了基于 Kahn 算法的拓扑排序程序，满足 ``文件读取 + 入度维护 + 队列处理 + 环检测'' 的要求，并提供了无环 / 有环的样例验证。程序结构清晰，错误处理完善，易于扩展与维护。拓扑排序在任务调度、依赖分析等领域有广泛应用，理解其原理与实现对数据结构学习具有重要意义。

\end{document}
